\chapter{Root Systems and Weyl Groups}
\label{ch:root-systems}

\section{Positive roots and chambers}

\begin{definition}[Finite root systems]\label{res:E01}
A reduced crystallographic root system in a Euclidean space \(E\)
is a finite spanning subset \(R\subset E\setminus\{0\}\) such that
\[
 s_\alpha(R)=R,\qquad
 \langle\beta,\alpha^\vee\rangle
 =\frac{2(\beta,\alpha)}{(\alpha,\alpha)}\in\mathbb Z,\qquad
 R\cap\mathbb R\alpha=\{\alpha,-\alpha\}.
\]
The coroots \(2\alpha/(\alpha,\alpha)\) form another root system,
the dual one. Orthogonal irreducible components are unique.
The rank two possibilities are \(A_1\sqcup A_1,A_2,B_2,G_2\),
with respectively \(4,6,8,12\) roots.
\end{definition}
\begin{proof}
The coroot reflection is the same orthogonal reflection, and
\[
 \frac{2(\alpha^\vee,\beta^\vee)}
      {(\alpha^\vee,\alpha^\vee)}
 =\frac{2(\alpha,\beta)}{(\beta,\beta)}.
\]
This proves the dual assertions and double duality.
For two independent roots making angle \(\theta\),
the product of the two integral Cartan numbers is
\(4\cos^2\theta\), a nonnegative integer less than \(4\).
The nonorthogonal cases give the squared length ratios
\(1,2,3\), and the acute angles \(60,45,30\) degrees.
To see that the rank two list is exhaustive, take two adjacent
root rays making the smallest positive angle. Reflection in
their perpendicular lines repeats this sector around the
circle. The angle restrictions give \(2,3,4,6\) pairs of
opposite rays. No extra root ray can occur inside a sector,
by its minimality; a second length on a ray is excluded by
reducedness. Integrality fixes equal lengths in \(A_2\)
and the alternating ratios \(2,3\) in \(B_2,G_2\).
The coordinate models below verify their existence.
Uniqueness of the orthogonal decomposition will also follow
from the simple-root basis proved next: it is the decomposition
of its graph into connected components.
\end{proof}

\begin{theorem}[Simple roots]\label{res:E02}
Choose \(v\in E\) off the root hyperplanes and put
\(R^+=\{\alpha:(v,\alpha)>0\}\). The positive roots that are
not sums of two positive roots form a basis
\(\Pi=\{\alpha_1,\ldots,\alpha_r\}\) of \(E\).
Every root is an integral combination of \(\Pi\), with all
coefficients nonnegative or all nonpositive.
The simple reflections \(s_i=s_{\alpha_i}\) generate the Weyl
group \(W=\langle s_\alpha:\alpha\in R\rangle\), and
\(R=W\Pi\).
\end{theorem}
\begin{proof}
If a positive root is decomposable, split it into two positive
roots and repeat. The values of \((v,-)\) on \(R^+\) form a
finite set of positive numbers, so the process terminates.
Thus every positive root is a sum of simple ones.
Distinct simple roots have nonpositive inner product: if
\((\alpha,\beta)>0\), inspection of the four rank two systems
gives \(\alpha-\beta\) or \(\beta-\alpha\) as a root, making
one of the two simple roots decomposable.
For linear independence, split a nonzero relation into
\(u=\sum_{i\in I}a_i\alpha_i=\sum_{j\in J}b_j\alpha_j\),
where \(I,J\) are disjoint and all coefficients positive.
Then \(u\ne0\), as pairing with \(v\) shows, whereas
\((u,u)=\sum a_ib_j(\alpha_i,\alpha_j)\le0\), a contradiction.
This proves the basis and the sign assertion.

If \(\beta=\sum n_i\alpha_i\) is positive and not simple,
some \((\beta,\alpha_i)>0\), since
\((\beta,\beta)=\sum n_i(\beta,\alpha_i)>0\).
The root \(s_i\beta=\beta-
\langle\beta,\alpha_i^\vee\rangle\alpha_i\)
has smaller height. It cannot be negative unless \(\beta\)
is a multiple of \(\alpha_i\), because all other coefficients
are unchanged. In the latter case reducedness makes it simple.
Thus repeated simple reflections take any root to a simple
root, proving \(R=W_0\Pi\) for \(W_0=\langle s_i\rangle\).
For \(\beta=w\alpha_i\), \(s_\beta=ws_iw^{-1}\in W_0\),
so \(W_0=W\).
One also obtains the useful lowering assertion:
if \(\beta>0\) is not simple, then \(\beta-\alpha_i\)
is a positive root for some \(i\). Indeed the root string
through \(\beta\) in the direction just chosen contains that
root. These string assertions in an abstract root system
can equivalently be checked in its rank two subsystem.
\end{proof}

\begin{theorem}[Chambers, length, and a fundamental domain]\label{res:E03}
The Weyl group acts simply transitively on the chambers in
\(E\setminus\bigcup_{\alpha\in R}\alpha^\perp\).
The chamber for \(\Pi\) is
\(C=\{x:(x,\alpha_i)>0\ \forall i\}\).
For \(w\in W\) its shortest simple-reflection word has length
\[
 \ell(w)=|\{\alpha>0:w\alpha<0\}|.
\]
Every \(W\)-orbit has exactly one point in \(\overline C\).
\end{theorem}
\begin{proof}
The chamber corresponding to a sign choice is an intersection
of open half spaces, hence convex. The positive-root
expansions show that the walls of \(C\) are precisely the
simple-root hyperplanes. A generic segment between two
chambers crosses one wall at a time. Reflection in that wall
takes one adjacent chamber to the other, so \(W\) is transitive.

We give the word argument establishing freeness, rather than
assuming it in the gallery argument. A simple reflection
permutes \(R^+\setminus\{\alpha_i\}\): coefficients outside
\(\alpha_i\) remain nonnegative, and only \(\alpha_i\)
can change sign. Suppose \(u=s_{i_1}\cdots s_{i_m}\)
and \(u\alpha<0\) for a positive root \(\alpha\).
Follow \(\alpha\) through this word from right to left.
At the first sign change the positive root must equal the
simple root being reflected. If the suffix before that step
is \(v\), then \(v\alpha=\alpha_{i_k}\) and
\(s_\alpha=v^{-1}s_{i_k}v\). Thus \(us_\alpha\)
has an expression obtained by deleting the \(k\)-th letter.
In a shortest expression \(w=us_i\), this implies
\(u\alpha_i>0\), for otherwise two letters could be deleted.
Counting the positive roots sent negative gives
\(N(us_i)=N(u)+1\) in that case. Induction shows that every
shortest word has length \(N(w)\).
If \(wC=C\), it sends no positive root negative, hence its
shortest word has length zero and \(w=1\).

Existence of a representative in \(\overline C\) follows
by continuity from transitivity on chambers, or by reflecting
successively across a wall on whose negative side the point lies.
For uniqueness we record a consequence of a reduced word
\(w=s_{i_1}\cdots s_{i_m}\). Each
\(\gamma_k=s_{i_1}\cdots s_{i_{k-1}}\alpha_{i_k}\)
is positive, by the preceding deletion argument. For dominant
\(\lambda\),
\[
 \lambda-w\lambda
 =\sum_{k=1}^m
   \langle\lambda,\alpha_{i_k}^\vee\rangle\gamma_k
 \quad\text{belongs to }\sum_i\mathbb R_{\ge0}\alpha_i.
\]
If \(\mu=w\lambda\) is also dominant, applying the same
assertion to \(w^{-1}\) and \(\mu\) puts \(\mu-\lambda\)
in that same pointed cone. Their sum is zero, so
\(\lambda=\mu\).
\end{proof}

\section{Dynkin classification}

\begin{proposition}[Cartan matrices and reconstruction]\label{res:E04}
Let \(a_{ij}=\langle\alpha_j,\alpha_i^\vee\rangle\).
Then \(a_{ii}=2\), \(a_{ij}\in\mathbb Z_{\le0}\) for
\(i\ne j\), and \(a_{ij}a_{ji}\in\{0,1,2,3\}\).
The matrix \(\operatorname{diag}((\alpha_i,\alpha_i))A\)
is symmetric positive definite. The Dynkin graph joins \(i,j\)
by \(a_{ij}a_{ji}\) edges, with its arrow pointing toward
the shorter root. It determines the root system up to
isometry and one positive scale on each irreducible component.
\end{proposition}
\begin{proof}
The angle and simple-root arguments already give the stated
entries and products. Symmetrization is twice the Gram matrix.
On a connected component the ratio of adjacent squared lengths
is \(a_{ji}/a_{ij}\), so the graph determines every length
relative to one chosen length. It therefore determines the
Gram matrix of a basis. The induced isometry identifies
the simple reflections, and Theorem~\ref{res:E02} then
identifies all roots as their orbits.
Orthogonal graph components give invariant orthogonal
subspaces and \(W\) is the product of their Weyl groups.
This also proves the uniqueness of the irreducible decomposition.
\end{proof}

\begin{theorem}[Exhaustion of finite Dynkin diagrams]\label{res:E05}
The connected finite reduced crystallographic diagrams are
\[
 A_n\ (n\ge1),\quad B_n,C_n\ (n\ge2),\quad D_n\ (n\ge4),
 \quad E_6,E_7,E_8,F_4,G_2,
\]
with \(B_2=C_2\) up to exchanging the two root lengths.
\end{theorem}
\begin{proof}
Normalize the simple roots to squared length \(2\).
Their symmetric Gram matrix \(M\) has diagonal entries \(2\)
and off-diagonal entries \(-\sqrt m\), where
\(m=a_{ij}a_{ji}\). Every principal submatrix is positive
definite. A cycle is impossible: on its vertices the vector
with all entries \(1\) has quadratic value at most zero.
A vertex with four neighbors is impossible: put \(2\) at
the center and \(1\) at its four neighbors; again the
quadratic value is at most zero.
Thus the graph is a tree of valence at most three.

There cannot be two trivalent vertices. Take the path joining
two consecutive ones and retain two neighboring leaves at
each end. Give all vertices on the path value \(2\) and
the four leaves value \(1\). For single edges its quadratic
value is zero; heavier edges can only decrease it.
A triple edge cannot have any attached vertex: the
three-vertex path with edge products \(3,1\) has determinant
\(8-2(3+1)=0\). Hence a triple edge gives just \(G_2\).

A trivalent vertex and a double edge cannot coexist.
Retain the path from that vertex to the nearest double edge,
two additional leaves at the trivalent end, and the vertex
just beyond the double edge. Eliminate the two leaves by
Schur complement: each subtracts \(1/2\) from the central
diagonal, leaving \(1\). Eliminating the last vertex across
the double edge subtracts \(2/2=1\) from the other end.
The remaining chain has endpoint diagonals \(1\), all interior
diagonals \(2\), and off-diagonals \(-1\), so its all-ones
vector has quadratic value zero. If the double edge is
incident to the trivalent vertex, these two endpoint
subtractions occur on the same diagonal, leaving the
one-by-one zero matrix instead. Two double edges are
excluded by the same endpoint elimination along
the path between them; if they share a vertex,
both subtractions again leave one zero diagonal.

It remains to classify simply laced trees with at most
one branch and chains with one double edge.
A simple chain is \(A_n\).
For a branch with arms of lengths \(a\le b\le c\),
the inverse endpoint entry of the length \(a\) chain
matrix is \(a/(a+1)\); this follows by the determinant
recurrence \(d_m=2d_{m-1}-d_{m-2}\), \(d_m=m+1\).
Its Schur complement at the branch is positive precisely when
\[
 \frac1{a+1}+\frac1{b+1}+\frac1{c+1}>1.
\]
Thus \(a=1\); if \(b=1\) one gets \(D_{c+3}\), and if
\(b=2\), one gets \(c=2,3,4\), namely \(E_6,E_7,E_8\).
No \(b\ge3\) is possible.
For a double edge separating chains of \(p,q\ge1\) vertices,
the determinant is
\[
 (p+1)(q+1)-2pq=2-(p-1)(q-1).
\]
Positivity gives \(p=1\) or \(q=1\), the diagrams \(B_n,C_n\),
or \(p=q=2\), the diagram \(F_4\).
The two orientations of \(F_4\) are interchanged by reversing
the chain. These cases exhaust all graphs and orientations.
Their actual root systems are constructed next, so the list
has neither missing cases nor merely formal candidates.
\end{proof}

\begin{theorem}[Coordinate constructions]\label{res:E06}
Every diagram in Theorem~\ref{res:E05} is realized by a root
system. The classical systems are
\[
\begin{array}{c|l|r}
\text{type}&\text{roots}&|R|\\\hline
A_n&\varepsilon_i-\varepsilon_j\text{ in }\sum x_i=0
 &(n+1)n\\
B_n&\pm\varepsilon_i,\ \pm\varepsilon_i\pm\varepsilon_j&2n^2\\
C_n&\pm2\varepsilon_i,\ \pm\varepsilon_i\pm\varepsilon_j&2n^2\\
D_n&\pm\varepsilon_i\pm\varepsilon_j&2n(n-1).
\end{array}
\]
The exceptional cardinalities are
\[
 |G_2|=12,\quad |F_4|=48,\quad
 |E_6|=72,\quad |E_7|=126,\quad |E_8|=240.
\]
\end{theorem}
\begin{proof}
Signed coordinate permutations verify the reflection axioms
for the classical lists. Their simple bases are
\(\varepsilon_i-\varepsilon_{i+1}\), with the final root
\(\varepsilon_n\) in \(B_n\), \(2\varepsilon_n\) in \(C_n\),
and with the two final roots
\(\varepsilon_{n-1}-\varepsilon_n,\varepsilon_{n-1}+\varepsilon_n\)
in \(D_n\). For \(A_n\) use \(n+1\) coordinates.
The dot products give the indicated graphs.

In \(\sum_{i=1}^3x_i=0\), take the \(G_2\) roots
\[
 \{\varepsilon_i-\varepsilon_j:i\ne j\}
 \ \cup\
 \{\pm(2\varepsilon_i-\varepsilon_j-\varepsilon_k):
                       \{i,j,k\}=\{1,2,3\}\}.
\]
A short reflection interchanges two coordinates.
The reflection in \(2\varepsilon_1-\varepsilon_2-\varepsilon_3\)
sends each short root to a short root and permutes the six
displayed long roots; inserting
\(x\mapsto x-(x,\alpha)\alpha/3\) verifies this directly.
The simple roots can be
\(\varepsilon_1-\varepsilon_2\) and
\(-2\varepsilon_1+\varepsilon_2+\varepsilon_3\),
with Cartan products \(3\).

In \(\mathbb R^4\), take
\[
 F_4=\{\pm\varepsilon_i\}\cup
 \{\pm\varepsilon_i\pm\varepsilon_j:i<j\}\cup
 \{\tfrac12(\pm\varepsilon_1\pm\varepsilon_2
                    \pm\varepsilon_3\pm\varepsilon_4)\}.
\]
The three parts have \(8,24,16\) members.
Signed permutations preserve the list. Reflection in a
half root preserves the long roots; on the short roots it
either negates that root, fixes it, or interchanges a
coordinate root with a half root. This follows by the
possible inner products \(0,\pm\tfrac12,\pm1\);
for long roots the nonzero inner product is \(\pm1\).
These computations prove both reflection invariance and
integrality. One simple basis is
\[
 \varepsilon_2-\varepsilon_3,\quad
 \varepsilon_3-\varepsilon_4,\quad
 \varepsilon_4,\quad
 \tfrac12(\varepsilon_1-\varepsilon_2-\varepsilon_3-\varepsilon_4).
\]
Its expansion coefficients for \(x\) are
\(x_1+x_2,\ 2x_1+x_2+x_3,\
3x_1+x_2+x_3+x_4,\ 2x_1\).
For every displayed root they are integers of one sign;
this also verifies directly that this is a simple basis.

For \(E_8\), let
\[
 L=\{z\in\mathbb Z^8:\textstyle\sum z_i\text{ even}\}
 \cup
 \{z+(\tfrac12,\ldots,\tfrac12):z\in\mathbb Z^8,\
                                      \textstyle\sum z_i\text{ even}\}.
\]
This is an integral even lattice. Its norm \(2\) vectors are
the \(112\) vectors \(\pm\varepsilon_i\pm\varepsilon_j\)
and the \(128\) half-coordinate vectors with an even
number of minus signs. For a norm \(2\) vector \(\alpha\),
the map \(x\mapsto x-(x,\alpha)\alpha\) preserves \(L\)
and norm. It therefore preserves precisely this root list.
Take the basis
\[
\begin{split}
 \alpha_1&=\tfrac12(\varepsilon_1+\varepsilon_8
                         -\varepsilon_2-\cdots-\varepsilon_7),\\
 \alpha_2&=\varepsilon_1+\varepsilon_2,\qquad
 \alpha_3=\varepsilon_2-\varepsilon_1,\\
 \alpha_4&=\varepsilon_3-\varepsilon_2,\quad
 \alpha_5=\varepsilon_4-\varepsilon_3,\quad
 \alpha_6=\varepsilon_5-\varepsilon_4,\quad
 \alpha_7=\varepsilon_6-\varepsilon_5,\quad
 \alpha_8=\varepsilon_7-\varepsilon_6 .
\end{split}
\]
Its graph has arms \(1,2,4\).
Here is an explicit way to check simplicity, without relying
only on a drawing. For \(x=\sum c_i\alpha_i\), recover
\[
\begin{aligned}
c_1&=2x_8,&c_8&=x_7+x_8,&c_7&=x_6+x_7+2x_8,\\
c_6&=x_5+x_6+x_7+3x_8,&
c_5&=x_4+x_5+x_6+x_7+4x_8,\\
c_4&=x_3+x_4+x_5+x_6+x_7+5x_8,&
c_2&=(x_1+x_2+c_4)/2,&c_3&=c_2+x_8-x_1 .
\end{aligned}
\]
For a half root the parity condition makes these integers,
all nonnegative when \(x_8=1/2\), and all nonpositive
when \(x_8=-1/2\). For an integral root, the same formulas
with its two nonzero coordinates give one sign throughout.
This proves simplicity and the \(E_8\) graph.

Intersect with \(x_7+x_8=0\) to obtain \(E_7\),
with basis \(\alpha_1,\ldots,\alpha_7\).
There are \(60+2\) integral roots and \(64\) half roots.
Intersect further with \(x_6=x_7=-x_8\) to obtain \(E_6\),
with basis \(\alpha_1,\ldots,\alpha_6\).
There are \(40\) integral roots and \(32\) half roots.
The intersections are root systems because the reflection
in a vector of the subspace preserves that subspace.
Their graphs have arms \(1,2,3\) and \(1,2,2\), respectively.
\end{proof}

\section{Lattices and useful dominant elements}

\begin{proposition}[Root and weight lattices]\label{res:E07}
Set
\[
 Q=\sum_i\mathbb Z\alpha_i,\quad
 Q^\vee=\sum_i\mathbb Z\alpha_i^\vee,\quad
 P=(Q^\vee)^*,\quad P^\vee=Q^*.
\]
Then \(Q\subset P\), \(Q^\vee\subset P^\vee\), and
the fundamental weights \(\omega_i\), characterized by
\(\langle\omega_i,\alpha_j^\vee\rangle=\delta_{ij}\),
form a basis of \(P\).
Dominant integral weights are precisely
\(\sum_i n_i\omega_i\), \(n_i\in\mathbb Z_{\ge0}\).
Each real orbit has a unique dominant representative.
Moreover \([P:Q]=\det A\).
\end{proposition}
\begin{proof}
The simple coroots are a basis of the coroot lattice:
the proof of Theorem~\ref{res:E02} applies to the dual
root system. Taking dual bases proves the assertion about
fundamental weights. Integral Cartan pairings put every root
in \(P\), and the dual argument gives the other inclusion.
The inequalities for the closed positive chamber are
exactly the nonnegativity of the fundamental-weight
coefficients. Orbit uniqueness is Theorem~\ref{res:E03}.
Finally the simple roots have coordinate columns \(A\)
in the fundamental-weight basis, so the index is the
absolute determinant, which is positive by symmetrization.
For example these determinants are \(n+1,2,2,4,3,2,1,1,1\)
for \(A_n,B_n,C_n,D_n,E_6,E_7,E_8,F_4,G_2\).
They can be computed by the chain recurrence and branch
Schur complements in Theorem~\ref{res:E05}.
\end{proof}

\begin{proposition}[The longest element and finite weight sets]\label{res:E08}
There is a unique \(w_0\in W\) carrying \(R^+\) to \(-R^+\).
It satisfies \(w_0^2=1\) and \(\ell(w_0)=|R^+|\).
The element
\[
 \rho=\tfrac12\sum_{\alpha>0}\alpha
\]
satisfies \(\langle\rho,\alpha_i^\vee\rangle=1\).
Similarly \(\rho^\vee=\frac12\sum_{\alpha>0}\alpha^\vee\)
satisfies \(\langle\alpha_i,\rho^\vee\rangle=1\).
For a finite-dimensional \(\mathfrak g\)-module its
\(\mathfrak h\)-weights are integral, their multiplicities
are \(W\)-invariant, and their set is saturated in the sense
that, if \(\mu\) occurs and \(n=\langle\mu,\alpha^\vee\rangle
\ge0\), then \(\mu-j\alpha\) occurs for \(0\le j\le n\).
\end{proposition}
\begin{proof}
The negative chamber is a chamber; simple transitivity gives
\(w_0\), and its square stabilizes the positive chamber.
The length formula proves maximality and the value stated.
A simple reflection permutes the positive roots other
than \(\alpha_i\), so \(s_i\rho=\rho-\alpha_i\).
Comparing with the reflection formula gives its pairing \(1\).
Apply the same argument to the dual system.

The action of \(\mathfrak h\) is diagonalizable on a finite
module by Theorem~\ref{res:C07}. Restriction to every root
triple and the \(\mathfrak{sl}_2\) classification give
integer coroot eigenvalues. The finite direct sum of
weight spaces in a coset \(\mu+\mathbb Z\alpha\)
decomposes into \(\mathfrak{sl}_2\) strings.
Each string is symmetric about zero and contains every
intermediate weight of the same parity. Thus the endpoints
are exchanged by \(s_\alpha\), with their multiplicities,
and a string containing the eigenvalue \(n\ge0\)
contains \(n,n-2,\ldots,-n\).
This proves both multiplicity invariance and saturation.
\end{proof}

\begin{theorem}[Highest roots and the principal triple]\label{res:E09}
For a simple \(\mathfrak g\), there is a unique highest root
\(\theta\), and every root is at most \(\theta\) in the
simple-root order. Choose simple triples \(e_i,f_i,h_i\).
If \(2\rho^\vee=\sum_i c_i h_i\), then
\[
 e=\sum_i e_i,\qquad h=2\rho^\vee,\qquad f=\sum_i c_i f_i
\]
form an \(\mathfrak{sl}_2\)-triple.
For \(m>0\), let \(r_m\) be the number of positive roots of
height \(m\), and let \(r_0=r\). Define the exponents by
giving \(m\) multiplicity \(r_m-r_{m+1}\). Then
\[
 \mathfrak g|_{\langle e,h,f\rangle}
       \simeq\bigoplus_{i=1}^r L_{2m_i},\qquad
 \sum_i m_i=|R^+|.
\]
Here \(L_a\) has highest weight \(a\) and dimension \(a+1\).
The Coxeter number used in this volume is
\(h_R=1+\operatorname{ht}\theta=1+\max m_i\).
\end{theorem}
\begin{proof}
Choose a maximal positive root \(\theta\).
Its root vector is killed by all positive root operators.
By PBW, the submodule of the adjoint module generated by
that vector is spanned by negative root monomials applied
to it: move positive factors right and Cartan factors
to the middle. All its weights are therefore at most
\(\theta\). Since the adjoint module of a simple algebra
is simple, this submodule is all of \(\mathfrak g\).
Thus every root is at most \(\theta\); a second maximal root
must equal it. This argument uses PBW, not the highest-weight
classification of Chapter~7.

Distinct simple positive and negative root spaces bracket
to zero, whereas \([e_i,f_i]=h_i\). Therefore
\([e,f]=h\). The identities
\(\alpha_i(2\rho^\vee)=2\) give the other two triple
relations. In the adjoint module the \(h\)-weights are
twice the root heights, with zero space \(\mathfrak h\).
All simple \(\mathfrak{sl}_2\)-summands consequently have
even highest weight. Each contributes one dimension to
zero weight and one to each positive even weight up to
its highest weight. Hence there are \(r\) summands and
the number of summands with highest weight exactly \(2m\)
is \(r_m-r_{m+1}\). Summing their contributions to positive
weights gives \(\sum m_i=|R^+|\).
There are no trivial summands since \(r_0=r_1=r\).
The top nonzero weight is \(2\operatorname{ht}\theta\),
proving the assertion about the largest exponent.
\end{proof}

For reference the complete height calculation is summarized here:
\[
\begin{array}{c|l|c}
 R&\{m_i\}&h_R\\ \hline
 A_n&1,2,\ldots,n&n+1\\
 B_n,C_n&1,3,\ldots,2n-1&2n\\
 D_n&1,3,\ldots,2n-3,\ n-1&2n-2\\
 G_2&1,5&6\\
 F_4&1,5,7,11&12\\
 E_6&1,4,5,7,8,11&12\\
 E_7&1,5,7,9,11,13,17&18\\
 E_8&1,7,11,13,17,19,23,29&30
\end{array}
\]
These entries can be checked without a representation formula:
use the coordinate lists of Theorem~\ref{res:E06}, expand each
positive root in the indicated basis, and form
\(\sum_{\alpha>0}t^{\operatorname{ht}\alpha}\).
For each row it equals
\(\sum_{m\text{ in the middle column}}(t+\cdots+t^m)\);
the explicit coefficient-recovery formulas there also cover
the exceptional systems. In type \(D_n\), the exponent \(n-1\)
has multiplicity two when \(n\) is even.
The definition of \(h_R\) here does not silently identify it
with the order of a product of simple reflections.

The diagram and lattice arguments follow
\cite[Sections 9--13]{humphreys} and
\cite[Sections 21--23 and 32]{etingof}.
